% Commenti speciali

	% !TEX encoding = UTF-8 Unicode
	% !TEX TS-program = pdflatex
	% !TEX spellcheck = it-IT
	% !TEX root = Dispensa/afed.tex


\thispagestyle{empty}

% Ringraziamenti 


Questa dispensa è stata scritta da Mattia Garatti e raccoglie gli appunti del corsi \emph{Analisi Funzionale} ed \emph{Equazioni Differenziali} tenuti dal Prof. Marco Squassina negli anni accademici 2023/2024 e 2024/2025 (con qualche aggiunta).

Realizzare una dispensa è un'operazione complessa, che richiede numerosi controlli. L'esperienza suggerisce che è praticamente impossibile realizzare un'opera priva di errori. Potete segnalare eventuali errori al seguente indirizzo di posta elettronica:
\[ \mail{mattiagaratti@gmail.com} \]

%\vspace{20 pt}

\noindent \textbf{Ringraziamenti} 
\begin{itemize}
	\item Si ringrazia il Professor Marco Degiovanni per le brillanti osservazioni.
	\item Si ringrazia il Dottor Marco Marzocchi per i consigli sull'utilizzo del linguaggio \LaTeX $\,$ e per la diffusione della passione per l'Analisi Matematica.
	\item Si ringrazia la Dottoressa Chiara Lonati per gli inserti di Fisica Matematica.
	\item Si ringrazia il Dottor Riccardo Moraschi per i disegni del Teorema di estensione, delle immersioni di Sobolev ed del Lemma di Hopf, nonché per le puntuali correzioni e i preziosi consigli.
	\item Si ringrazia il Dottor Gabriele Geroldi per l'aiuto nella revisione di bozze.
	\item Si ringrazia il Dottor Luca Tamanini per le fruttuose discussioni e la revisione di bozze per le parti riguardanti gli spazi riflessivi e la rappresentazione degli elementi del duale di uno spazio di Lebesgue.
\end{itemize}

%\vspace{20 pt}

\noindent \textbf{Nota dell'autore} Non tutte le $C$ sono uguali.

\vfill

\begin{flushright}
	Ultimo aggiornamento: \today
\end{flushright}